%!TEX root = ../main_thesis.tex

\chapter{Linearizzazione e forma standard}
\label{cap:linearizzazione}

Il \texttt{Model} prodotto dall'analisi semantica (Capitolo~\ref{cap:semantica})
è una descrizione concreta del problema, ma può ancora contenere costrutti non
direttamente trattabili da un solver lineare. Questo capitolo descrive il
\emph{back-end} del compilatore: la traduzione del modello in una forma lineare
e la sua successiva riduzione alla forma standard, prerequisito del metodo del
simplesso.

\section{Concetti teorici}
\label{sec:linearizzazione-teoria}

\paragraph{Linearità.}
Un'espressione è \emph{lineare} nelle variabili decisionali se è una
combinazione lineare di queste, cioè una somma di termini ciascuno costituito da
una variabile moltiplicata per un coefficiente costante. Sono lineari le somme e
le differenze di espressioni lineari e il prodotto di una costante per una
variabile; non lo è, invece, il prodotto di due variabili o la divisione per
un'espressione che contenga variabili. Poiché i solver di programmazione lineare
operano per definizione su modelli lineari, ogni costrutto non lineare deve
essere eliminato o riformulato prima della risoluzione.

\paragraph{Riformulazione di minimo e massimo.}
Alcuni costrutti non lineari possono essere resi lineari introducendo
\emph{variabili ausiliarie} e vincoli aggiuntivi. È il caso del minimo e del
massimo di un insieme di espressioni. Per rappresentare $t = \min(e_1, \dots,
e_k)$ si introduce una nuova variabile $t$ e si impone

\begin{equation}
\label{eq:min-lin}
t \le e_i \qquad \forall i \in \{1, \dots, k\},
\end{equation}

sostituendo l'occorrenza del minimo con $t$. Dualmente, per il massimo si
impone $t \ge e_i$ per ogni $i$. Questa riformulazione è \emph{unilaterale}:
vincola $t$ da un solo lato ed è quindi esatta soltanto quando è
l'ottimizzazione stessa a spingere $t$ contro il valore vero, come accade per
un massimo che va minimizzato o per un minimo che va massimizzato.

\paragraph{Il metodo Big-M.}
Negli altri casi occorre imporre anche il vincolo dal lato opposto, che però
deve valere solo per uno degli argomenti: quello che realizza davvero il
minimo o il massimo. La tecnica classica per esprimere una scelta di questo
tipo è il metodo \emph{Big-M}: alla disuguaglianza si somma un termine
$M(1-\delta)$, dove $\delta$ è una variabile binaria e $M$ una costante
sufficientemente grande. Quando $\delta = 1$ il termine si annulla e il
vincolo è attivo; quando $\delta = 0$ la costante $M$ rende la disuguaglianza
sempre soddisfatta, disattivandola di fatto. Un valore di $M$ più grande del
necessario peggiora però la qualità numerica del modello, quindi conviene
ricavare da estremi noti il più piccolo $M$ sufficiente. Con intervalli
$e_i \in [L_i, U_i]$ e $U_t = \max_i U_i$, il valore $t = \max(e_1, \dots,
e_k)$ è descritto esattamente da

\begin{equation}
\label{eq:max-exact}
\begin{aligned}
t &\ge e_i && \forall i,\\
t &\le e_i + (U_t-L_i)(1-\delta_i) && \forall i,\\
\textstyle\sum_i \delta_i &= 1.
\end{aligned}
\end{equation}

La prima riga impone che $t$ non sia inferiore ad alcun argomento; la seconda
è il vincolo opposto, attivo solo per l'argomento scelto da $\delta_i$, con
$(U_t-L_i)$ nel ruolo di $M$; la terza richiede che esattamente un argomento
sia scelto. Il minimo si tratta in modo simmetrico, scambiando i versi delle
disuguaglianze.

\paragraph{Riformulazione del valore assoluto.}
Il valore assoluto si riconduce al caso precedente osservando che $|e| =
\max(e, -e)$: per rappresentare $t = |e|$ si impone

\begin{equation}
\label{eq:abs-lin}
t \ge e, \qquad t \ge -e,
\end{equation}

con la stessa avvertenza sulla direzione di ottimizzazione; poiché $|e| \ge
0$, la variabile ausiliaria può inoltre essere dichiarata non negativa. Se
l'intervallo di $e$ attraversa lo zero e serve il valore esatto, un selettore
binario disattiva uno dei due rami secondo lo schema Big-M appena visto. Se
invece il segno di $e$ è già noto, il valore assoluto coincide direttamente
con $e$ o con $-e$, senza variabili ausiliarie.

\paragraph{La forma standard.}
Il metodo del simplesso richiede che il problema sia espresso in \emph{forma
standard}: funzione obiettivo di minimizzazione, vincoli di sola uguaglianza e
variabili tutte non negative. La riduzione a questa forma si ottiene con
trasformazioni meccaniche. Un vincolo di disuguaglianza viene reso
un'uguaglianza aggiungendo una variabile non negativa: una variabile di
\emph{slack} per i vincoli di tipo $\le$ e una variabile di \emph{surplus} per
quelli di tipo $\ge$,

\begin{equation}
\label{eq:slack}
\textstyle\sum_i a_i x_i \le b \;\Longrightarrow\; \sum_i a_i x_i + s = b, \quad s \ge 0.
\end{equation}

Una variabile \emph{libera} $x$, che può assumere valori negativi, viene
sostituita dalla differenza di due variabili non negative, $x = x' - x''$ con
$x', x'' \ge 0$. Infine un problema di massimizzazione viene ricondotto a uno di
minimizzazione osservando che $\max\, c^\top x = -\min\, (-c^\top x)$.

\section{Applicazione in ROOC: da \texttt{Model} a \texttt{LinearModel}}
\label{sec:linear-model}

La linearizzazione percorre ricorsivamente l'albero di ogni espressione e ne
costruisce una rappresentazione lineare: una mappa che associa a ogni
variabile il proprio coefficiente, più un termine costante. Addizione e
sottrazione combinano le mappe termine a termine, la negazione moltiplica i
coefficienti per $-1$, mentre moltiplicazione e divisione sono ammesse solo
quando l'altro operando è costante. Un prodotto tra due espressioni che
contengono variabili è intrinsecamente non lineare e viene segnalato come
errore.

\paragraph{Analisi degli intervalli.}
Prima di generare i vincoli, il compilatore calcola per ogni variabile un
intervallo conservativo, partendo dagli estremi dichiarati nel dominio e
raffinandoli con le informazioni contenute nei vincoli del modello. Questi
intervalli forniscono i valori $L_i$ e $U_i$ richiesti dalle formulazioni
esatte della Sezione~\ref{sec:linearizzazione-teoria}, senza ricorrere a
costanti Big-M arbitrarie.

\paragraph{Direzione d'uso dei valori.}
Durante la visita, il compilatore tiene inoltre traccia di come ogni valore
concorre all'ottimo: se l'ottimizzazione lo spinge verso il basso, verso
l'alto, oppure se il valore deve essere esatto, come accade in un vincolo di
uguaglianza. È questa informazione a decidere, per i costrutti non lineari,
se basta la riformulazione unilaterale o se serve quella esatta. 
Il risultato è che ogni vincolo viene portato nella forma di una combinazione lineare confrontata
con una costante e l'obiettivo in un vettore di coefficienti. Il risultato è
un \texttt{LinearModel}: l'elenco delle variabili effettivamente usate, i loro
domini, l'obiettivo e i vincoli lineari (\texttt{LinearConstraint}).


\section{Linearizzazione di minimo, massimo e valore assoluto}
\label{sec:linearizzazione-minmax}

Per minimo e massimo la formulazione è scelta in base alla direzione d'uso.
Quando questa rende sufficiente la riformulazione unilaterale, il compilatore
introduce una sola variabile reale e i vincoli \eqref{eq:min-lin}, o i
corrispondenti per il massimo: non servono variabili binarie e un problema
lineare continuo resta tale. Quando invece è richiesto il valore esatto, viene
applicata la formulazione \eqref{eq:max-exact}, o la sua simmetrica per il
minimo, con i coefficienti Big-M ricavati dagli intervalli calcolati. Se gli
estremi finiti necessari non sono disponibili, la compilazione termina con un
errore esplicito che indica le variabili prive di limiti, anziché introdurre
una costante arbitraria.

Il valore assoluto segue lo stesso criterio: se il segno dell'argomento è noto
viene sostituito direttamente da $e$ o da $-e$; altrimenti si usa il solo
epigrafo \eqref{eq:abs-lin} quando la direzione d'uso lo consente, e il
selettore binario negli altri casi.

\section{Linearizzazione degli operatori logici}
\label{sec:linearizzazione-logica}

Gli operatori logici del linguaggio (Capitolo~\ref{cap:linguaggio}) vengono
anch'essi eliminati in questa fase. Quando un'espressione logica è usata come
valore numerico, ad esempio nell'obiettivo, viene applicata una
\emph{reificazione}: la sottoespressione è sostituita da una variabile
binaria (\texttt{\$and\_0}, \texttt{\$or\_0}, e così via), legata ai propri
operandi da vincoli lineari. Poiché il sistema dei tipi garantisce valori pari
a $0$ o $1$, sono applicabili le codifiche classiche della programmazione
lineare intera. Per la congiunzione $z = e_1 \land \dots \land e_k$ si impone

\begin{equation}
\label{eq:and-lin}
z \le e_i \quad \forall i, \qquad z \ge \textstyle\sum_i e_i - (k - 1),
\end{equation}

mentre per la disgiunzione $z = e_1 \lor \dots \lor e_k$ si impone

\begin{equation}
\label{eq:or-lin}
z \ge e_i \quad \forall i, \qquad z \le \textstyle\sum_i e_i.
\end{equation}

La negazione non richiede alcuna variabile ausiliaria, perché $\lnot e$ è
l'espressione affine $1 - e$. L'implicazione $z = (a \Rightarrow b)$, essendo
equivalente a $\lnot a \lor b$, si codifica con $z \ge 1 - a$, $z \ge b$ e $z
\le 1 - a + b$; la doppia implicazione e la disgiunzione esclusiva sono
trattate in modo analogo con quattro disuguaglianze ciascuna, mentre lo
\texttt{xor} su più di due operandi è ridotto a una catena di \texttt{xor}
binari.

Un vincolo scritto come sola espressione logica è un'asserzione: la formula
deve risultare vera, ma questo non richiede necessariamente la reificazione
della radice.
Se gli operandi sono affini, l'asserzione viene tradotta direttamente: per
$e_1 \land \dots \land e_k$ vero si impone $\sum_i e_i = k$, per la
disgiunzione vera si impone $\sum_i e_i \ge 1$, e per $a \Rightarrow b$ si
impone $a \le b$. Le forme negate sono ottenute in modo analogo.

Le formule annidate vengono invece visitate secondo il valore che devono
certificare. Una congiunzione richiesta vera genera le condizioni sui singoli
figli; una disgiunzione richiesta falsa procede nello stesso modo sui valori
falsi. Nei casi disgiuntivi vengono introdotti soltanto \emph{testimoni}
binari unidirezionali per i rami che possono certificare l'asserzione. Tali
testimoni non sono mai usati come valori aritmetici e non devono quindi
rappresentare l'intera formula in entrambe le direzioni. Per esempio, il
vincolo di copertura $x_u \lor x_v$ diventa direttamente
$x_u+x_v \ge 1$, senza una variabile \texttt{\$or\_N}.

La reificazione esatta rimane invece necessaria quando il valore logico entra
in un'espressione aritmetica: in quel caso la variabile ausiliaria coincide con
il valore di verità in ogni soluzione ammissibile. In entrambi i casi la
riformulazione avviene soltanto nel passaggio a \texttt{LinearModel}; il
\texttt{Model} conserva la struttura logica originale, disponibile per
eventuali risolutori futuri in grado di trattarla nativamente.

\section{Trasformazione in forma standard}
\label{sec:forma-standard}

L'ultima trasformazione riduce il \texttt{LinearModel} a uno
\texttt{StandardLinearModel}, applicando i passaggi teorici visti in precedenza.

\paragraph{Vincolo sulle variabili.}
La forma standard, così come è utilizzata dal simplesso, è definita per
variabili \emph{continue}. Per questo motivo la trasformazione richiede che
tutte le variabili siano di tipo reale (con o senza vincolo di non negatività);
in presenza di variabili intere o binarie restituisce un errore di dominio non
valido. La risoluzione di modelli con variabili intere segue infatti un percorso
diverso, descritto nel Capitolo~\ref{cap:solver}.

\paragraph{I passaggi della trasformazione.}
La riduzione procede in quattro fasi. Anzitutto, per ogni variabile dotata di
estremi finiti vengono aggiunti i vincoli espliciti che li impongono (un vincolo
$\ge$ per il limite inferiore e un vincolo $\le$ per quello superiore). In
secondo luogo, ogni variabile libera viene sostituita dalla differenza di due
variabili non negative, riportando i relativi coefficienti nei vincoli e
nell'obiettivo. In terzo luogo, ogni vincolo viene normalizzato a un'uguaglianza
secondo la~\eqref{eq:slack}, aggiungendo una variabile di slack o di surplus a
seconda del verso della disuguaglianza. Infine, se il problema è di
massimizzazione, i coefficienti dell'obiettivo vengono negati e viene registrato
un apposito segnale, così da ricondurlo a una minimizzazione.

\paragraph{Il modello in forma standard.}
Il risultato è uno \texttt{StandardLinearModel},
che contiene l'elenco (ordinato) delle variabili, i coefficienti dell'obiettivo,
i vincoli ridotti a uguaglianze, il termine costante dell'obiettivo
(\texttt{objective\_offset}) e un indicatore (\texttt{flip\_objective}) che
ricorda se l'obiettivo originale era di massimizzazione. Questi ultimi due valori
permettono, una volta risolto il problema, di riconvertire il valore ottimo
trovato in quello del problema di partenza.

A questo punto il modello è pronto per essere tradotto nel \emph{tableau} del
simplesso e risolto, come si vedrà nel prossimo capitolo.
